%% Signe du travail d'une force selon l'angle alpha entre F et le déplacement AB
%% (leçon pc-travail-energie) : W = F x AB x cos(alpha).
%% (1) alpha = 40° : moteur — (2) alpha = 90° : nul — (3) alpha = 140° : résistant.
\documentclass[tikz,border=4pt]{standalone}
\usepackage{liaisons}
\begin{document}
\begin{tikzpicture}[x=1cm,y=1cm,
  vect/.style={-{Stealth[length=6pt]}, line width=1.2pt},
  croix/.style={line width=0.9pt, black}]

%% --- macro locale : un schéma complet, paramétré par l'angle alpha ----
%%  #1 décalage horizontal, #2 numéro, #3 angle alpha, #4 légende
\newcommand\cas[4]{%
  \begin{scope}[shift={(#1,0)}]
    \coordinate (A) at (0,0);
    \coordinate (B) at (2.2,0);
    % secteur de l'angle alpha (en fond)
    \fill[solideE!25] (A) -- (0.55,0) arc (0:#3:0.55) -- cycle;
    \draw[line width=0.5pt, black!55] (0.55,0) arc (0:#3:0.55);
    \node[font=\small] at ({#3/2}:0.82) {$\alpha$};
    % le déplacement AB
    \draw[vect, solideB] (A) -- (B);
    \node[font=\small, text=solideB, anchor=north, inner sep=3pt] at (1.1,-0.06)
         {$\overrightarrow{\mathrm{AB}}$};
    % les points A et B (croix)
    \foreach \px in {0,2.2} {
      \draw[croix] (\px-0.09,-0.09) -- (\px+0.09,0.09)
                   (\px-0.09,0.09) -- (\px+0.09,-0.09); }
    \node[font=\small, anchor=north east, inner sep=2pt] at (A) {A};
    \node[font=\small, anchor=north west, inner sep=2pt] at (B) {B};
    % la force F, appliquée en A
    \draw[vect, solideA] (A) -- (#3:1.5)
         node[anchor=south, inner sep=2pt, font=\small, text=solideA] {$\vec{F}$};
    % numéro et légende
    \node[font=\small, anchor=south] at (1.1,1.75) {\textbf{(#2)}};
    \node[font=\footnotesize, anchor=north, align=center] at (1.1,-0.75) {#4};
  \end{scope}}

\cas{0}{1}{40}{$\alpha < 90^\circ$\\ travail \textbf{moteur}\\ $W > 0$}
\cas{3.6}{2}{90}{$\alpha = 90^\circ$\\ la force \textbf{ne travaille pas}\\ $W = 0$}
\cas{7.2}{3}{140}{$\alpha > 90^\circ$\\ travail \textbf{résistant}\\ $W < 0$}
\end{tikzpicture}
\end{document}
